\clearpage
\VerdierPhysicalPageAnchor{0165}%
\begingroup\small\setlength{\parskip}{0pt}\setlength{\emergencystretch}{1.5em}
\VerdierSemanticAnchor{ch4-4-2-1-inclusions}
こうして $\mathcal{F}\ell_2(\mathcal J)$ から $\mathcal A$ への二つの
関手を定義した。定義により、次の包含関係がある：
\[
B\!\left[\begin{smallmatrix}n\\f,g\end{smallmatrix}\right]\subset H^n(f),
\qquad
Z\!\left[\begin{smallmatrix}n\\f,g\end{smallmatrix}\right]\subset H^n(g).
\]
$\omega_3=(\mathord\bullet\xrightarrow f\mathord\bullet\xrightarrow g
\mathord\bullet\xrightarrow h\mathord\bullet)$ を
$\mathcal{F}\ell_3(\mathcal J)$ の対象とする。

\VerdierSemanticAnchor{ch4-4-2-2}
\medskip\noindent\textbf{命題 4.2.2.}\quad
\textit{次の包含関係がある：}
\[
B\!\left[\begin{smallmatrix}n\\g,h\end{smallmatrix}\right]
\subset
Z\!\left[\begin{smallmatrix}n\\f,g\end{smallmatrix}\right]
\subset H^n(g).
\]
第二の包含関係は参照のために記した。第一の包含関係を証明するには、次の列の
二つの射の合成が零であることを示せば十分である：
\VerdierSemanticAnchor{ch4-4-2-2-proof-sequence}
\[
H^{n-1}(h)\xrightarrow{\delta_1}H^n(g)
\xrightarrow{\delta_2}H^{n+1}(f)
\]
ところが、対象
$(\mathord\bullet\xrightarrow g\mathord\bullet\xrightarrow h\mathord\bullet)$
は $\mathcal{F}\ell_2(\mathcal J)$ の対象であり、完全列
\[
H^{n-1}(h)\xrightarrow{\delta_1}H^n(g)
\xrightarrow u H^n(hg),
\]
を与え、$\mathcal{F}\ell_2(\mathcal J)$ の射
\VerdierSemanticAnchor{ch4-4-2-2-fl2-morphism}
\[
\begin{tikzcd}[column sep=4em,row sep=2.2em]
{}\arrow[r,"f"]\arrow[d,"\operatorname{id}"']&{}
 \arrow[r,"g"]\arrow[d,"\operatorname{id}"']&{}\arrow[d,"h"]\\
{}\arrow[r,"f"']&{}\arrow[r,"hg"']&{}
\end{tikzcd}
\]
は可換図式
\VerdierSemanticAnchor{ch4-4-2-2-commutative-triangle}
\[
\begin{tikzcd}[column sep=5em,row sep=2.7em]
H^n(g)\arrow[r,"\delta_2"]\arrow[d,"u"']&H^{n+1}(f)\\
H^n(hg)\arrow[ur,"\delta_3"']&{}
\end{tikzcd}
\]
を与える。したがって命題が従う。
\endgroup
